\newpage
\phantomsection
\label{prf:capstone-functions}

\begin{remark*}[Return]
\hyperref[chap:functions]{Return to Chapter: Functions}
\end{remark*}

\subsection*{Capstone --- Functions}

\begin{remark*}[Capstone Problem]
Let $A \subseteq \mathbb{R}$, let $c$ be a cluster point of $A$, and let
$f, g : A \to \mathbb{R}$. Throughout, assume all stated limits exist and
are finite.

\begin{enumerate}[label=(\alph*)]

\item \textbf{(Uniqueness via the Sequential Criterion.)}
Suppose $\lim_{x \to c} f(x) = L$ and $\lim_{x \to c} f(x) = L'$. Without
using the Uniqueness of Limits theorem directly, prove $L = L'$ by applying
the Sequential Criterion: construct a sequence $(x_n) \subseteq A \setminus \{c\}$
with $x_n \to c$ (using the fact that $c$ is a cluster point), and use
the uniqueness of sequence limits to conclude $L = L'$.

\item \textbf{(Algebra via the bridge principle.)}
Suppose $\lim_{x \to c} f(x) = L$ and $\lim_{x \to c} g(x) = M$.
Prove the Product Rule $\lim_{x \to c}(fg)(x) = LM$ by the following
route: take any $(x_n) \subseteq A \setminus \{c\}$ with $x_n \to c$,
apply the Sequential Criterion to get $f(x_n) \to L$ and $g(x_n) \to M$,
then invoke the Product Rule for sequence limits to get $f(x_n)g(x_n) \to LM$,
and conclude via the Sequential Criterion.

\item \textbf{(Local boundedness as a prerequisite.)}
Before proving the Product Rule in full generality, the Limit Implies Local
Boundedness theorem is needed. State precisely what it asserts (as a quantified
statement), and explain at which step of part (b) it would be invoked if you
were instead proving the Product Rule directly from the $\varepsilon$-$\delta$
definition without using the Sequential Criterion.

\item \textbf{(Heine--Borel in the limit context.)}
Let $X \subseteq \mathbb{R}$ be closed and bounded, and let $h : X \to \mathbb{R}$.
Suppose $(x_n) \subseteq X$ is a sequence with $x_n \to c$ for some $c \in \mathbb{R}$.

\begin{enumerate}[label=(\roman*)]
\item Use Heine--Borel to prove $c \in X$.
\item Deduce that if $\lim_{x \to c;\, x \in X} h(x) = L$ exists, then $c$ is
a cluster point of $X$ (justify this), and the Sequential Criterion applies.
\item Explain why the assumption that $X$ is closed is essential in (i): give
an explicit example of a bounded but not closed set $X$ and a sequence $(x_n)
\subseteq X$ with $x_n \to c \notin X$.
\end{enumerate}

\item \textbf{(Counterexample: composition of limits without the continuity
hypothesis.)}
The Composition of Limits theorem requires the extra hypothesis $g(c_2) = L$
when $c_2 \in B$. Exhibit explicit functions $f$ and $g$ and a point $c_1$
such that $\lim_{x \to c_1} f(x) = c_2$ and $\lim_{y \to c_2} g(y) = L$ both
exist, but $\lim_{x \to c_1} g(f(x)) \neq L$. Identify exactly which hypothesis
fails in your example.

\end{enumerate}

\medskip
\textit{Concepts synthesized:}
Cluster points (sequential characterisation),
Sequential Criterion for function limits,
Uniqueness of limits,
Limit Implies Local Boundedness,
Product Rule for function limits,
Heine--Borel,
Closed iff Sequentially Closed,
Composition of Limits (failure mode).

\medskip
\textit{Connection to simulation.}
The bridge principle in part (b) is structurally the same operation as
lifting a numerical scheme's convergence proof from scalar sequences to
function-valued processes. In the Brownian motion simulation, the drift
coefficient $b(\theta)$ and diffusion coefficient $\sigma(\theta)$ on the
torus both depend continuously on position $\theta$; the product
$b(\theta)\sigma(\theta)$ appearing in the It\^{o} formula inherits its
limit behaviour at boundary identifications via exactly this transfer
mechanism. Part (d) is the compactness argument that guarantees the torus
$\mathbb{T}^2$ (as a quotient of the closed square $[0,1]^2$) is a valid
state space: closed and bounded, so every trajectory has a convergent
subsequence with its limit still on the torus.
\end{remark*}

\begin{proof}
\textbf{Solution.}~\\

\textbf{Part (a): Uniqueness via the Sequential Criterion.}

Since $c$ is a cluster point of $A$, for each $n \in \mathbb{N}$ there exists
$x_n \in A \setminus \{c\}$ with $|x_n - c| < 1/n$. The sequence $(x_n) \subseteq
A \setminus \{c\}$ satisfies $x_n \to c$. Applying the Sequential Criterion to
$\lim_{x \to c} f(x) = L$ gives $f(x_n) \to L$. Applying it to
$\lim_{x \to c} f(x) = L'$ gives $f(x_n) \to L'$. Since a sequence has at most
one limit, $L = L'$.

\medskip
\textbf{Part (b): Product Rule via the bridge principle.}

Let $(x_n) \subseteq A \setminus \{c\}$ with $x_n \to c$ be arbitrary. By the
Sequential Criterion applied to $\lim_{x \to c} f(x) = L$: $f(x_n) \to L$. By the
Sequential Criterion applied to $\lim_{x \to c} g(x) = M$: $g(x_n) \to M$. By
the Product Rule for sequence limits: $f(x_n)g(x_n) \to LM$. Since $(x_n)$ was
arbitrary, the Sequential Criterion (reverse direction) gives
$\lim_{x \to c}(fg)(x) = LM$.

\medskip
\textbf{Part (c): Role of Local Boundedness.}

The Limit Implies Local Boundedness theorem asserts:
\[
  \lim_{x \to c} f(x) = L
  \;\Rightarrow\;
  \exists \delta > 0\;\exists M > 0\;\forall x \in A\;
  \bigl(0 < |x-c| < \delta \Rightarrow |f(x)| \leq M\bigr).
\]
In a direct $\varepsilon$-$\delta$ proof of the Product Rule, one writes:
\[
  |f(x)g(x) - LM|
  \leq |f(x)||g(x) - M| + |M||f(x) - L|.
\]
The term $|f(x)|$ must be bounded by a constant before the $|g(x) - M|$ factor
can be made small. Local Boundedness supplies this bound: it gives a $\delta_0$
and $M_0$ such that $|f(x)| \leq M_0$ for $x$ in the punctured $\delta_0$-ball.
Without this bound, $|f(x)|$ could grow as $x \to c$ and prevent the product
from converging. The Sequential Criterion route in part (b) sidesteps this issue
because the sequence $f(x_n)$ is already convergent (hence bounded by the
analogous sequence result), so the product rule for sequences handles the
boundedness implicitly.

\medskip
\textbf{Part (d): Heine--Borel in the limit context.}

\textit{(i)} The sequence $(x_n) \subseteq X$ is bounded (since $X$ is bounded)
and converges to $c$. By Heine--Borel, $X$ is closed. Since $X$ is closed and
$(x_n) \subseteq X$ converges to $c$, the Closed iff Sequentially Closed
characterisation gives $c \in X$.

\textit{(ii)} Since $X$ is a non-degenerate subset of $\mathbb{R}$ and the sequence
$(x_n)$ witnesses that $c$ can be approximated by distinct points of $X$
(if $x_n \neq c$ for infinitely many $n$, those terms form a subsequence in
$X \setminus \{c\}$ converging to $c$, so $c$ is a cluster point of $X$).
The Sequential Criterion then applies: $\lim_{x \to c;\, x \in X} h(x) = L$
is equivalent to $h(x_n) \to L$ for every such sequence.

\textit{(iii)} Take $X = (0,1]$ (bounded, not closed) and $x_n = 1/n \in X$.
Then $x_n \to 0 \notin X$. A function $h$ defined on $(0,1]$ could have
$\lim_{x \to 0^+} h(x) = L$ without $h$ being defined or extendable at $0$;
the limit point escapes the domain.

\medskip
\textbf{Part (e): Composition failure without the continuity hypothesis.}

Define $f : \mathbb{R} \to \mathbb{R}$ by $f(x) = 0$ for all $x$, and define
$g : \mathbb{R} \to \mathbb{R}$ by
\[
  g(y) = \begin{cases} 1 & \text{if } y = 0, \\ 0 & \text{if } y \neq 0. \end{cases}
\]
Set $c_1 = 1$ and $c_2 = 0$. Then $\lim_{x \to 1} f(x) = 0 = c_2$ and
$\lim_{y \to 0} g(y) = 0 = L$. However, $g(f(x)) = g(0) = 1$ for all $x$, so
$\lim_{x \to 1} g(f(x)) = 1 \neq 0 = L$.

The failing hypothesis is $g(c_2) = L$: here $g(0) = 1 \neq 0 = L$. The function
$f$ sends every input to $c_2 = 0$, so $f(x)$ always lands on the single point
where $g$ disagrees with its limit. Since the punctured condition on $g$ excludes
$y = c_2$ from consideration, the limit $\lim_{y \to 0} g(y) = 0$ never ``sees''
the bad value $g(0) = 1$; but the composition forces $g$ to evaluate at $0$
constantly.
\end{proof}

\begin{remark*}[Dependencies]~\\
\begin{itemize}
  \item \hyperref[def:cluster-point-r]{Definition: Cluster Point}
  \item \hyperref[thm:cluster-point-sequential]{Theorem: Sequential Characterization of Cluster Points}
  \item \hyperref[prop:sequential-criterion-limits]{Proposition: Sequential Criterion for Function Limits}
  \item \hyperref[thm:limit-unique]{Theorem: Uniqueness of Limits of Functions}
  \item \hyperref[thm:limit-implies-local-bounding]{Theorem: Limit Implies Local Boundedness}
  \item \hyperref[thm:product-rule-function-limits]{Theorem: Product Rule for Function Limits}
  \item \hyperref[thm:heine-borel]{Theorem: Heine--Borel Theorem for $\mathbb{R}$}
  \item \hyperref[cor:closed-iff-seq-limits]{Corollary: Closed iff Sequentially Closed}
  \item \hyperref[thm:composition-of-limits]{Theorem: Composition of Limits}
\end{itemize}
\end{remark*}
